\input{../header}
\usepackage{pgfplots}
\rhead{Your name: \rule{8cm}{0.15mm}}

\begin{document}

\section{Learning target DFa, version 1}

Correctly interpret the meaning of the derivative in context. Assign the correct units to the value of the derivative.

\hrulefill

A company manufactures rope, and the total cost of producing $r$ feet of rope is $C(r)$ dollars.
\begin{enumerate}[leftmargin=0pt]
    \item Suppose $C(2000) = 800$. What are the units on the 2000? What are the units on the 800?
    \vfill
    \item Write a sentence explaining what $C(2000) = 800$ means in the story.
    \vfill
    \item Suppose $C'(2000) = 0.35$. What are the units on the 2000? What are the units on the 0.35?
    \vfill
    \item Write a sentence explaining what $C'(2000) = 0.35$ means in the story. \\
    \textbf{Don't say ``per,'' and don't say ``rate.''}
    \vfill
\end{enumerate}


%%%%%%%%%%%%%%%%%%%%%%%%%%%%%%%%%%%%%%%%%%%%%%%%%%%%%%%%%
\pagebreak
%%%%%%%%%%%%%%%%%%%%%%%%%%%%%%%%%%%%%%%%%%%%%%%%%%%%%%%%%

\section{Learning targets DF1 and DFa, version 2}

Correctly interpret the meaning of the derivative in context. Assign the correct units to the value of the derivative.

\hrulefill

Arapaho Glacier is a mountain glacier in Roosevelt National Forest,
west of Boulder, CO. The following table\footnote{
Haugen, B., Scambos, T., Pfeffer, T., \& Anderson, R. (2010). Twentieth-century changes in the thickness and extent of Arapaho Glacier, Front Range, Colorado. \textit{Arctic, Antarctic, and Alpine Research, 42}(2), 198-209.}
gives the surface area, $A(t)$, in square meters, of Arapaho Glacier in the year $t$. 
\begin{center}
	\begin{tabular}{|c|c|c|c|c|}
	\hline
	$t$ & 1900 & 1960 & 1973 & 1999\\
	\hline
	$A(t)$ & 338,282 & 250,764 & 225,000 & 162,027\\
	\hline
	\end{tabular}
 \end{center}
	
\begin{enumerate}[leftmargin=0pt]

\item Compute an approximation for $A'(1960)$, and {\bf include units} for this number.
\vfill
\item Write a sentence explaining what the number means about how the area of the glacier is changing.\\
\textbf{Don't say ``per,'' and don't say ``rate.''}

\vfill
\vfill

\item Compute an approximation for $A'(1999)$, and {\bf include units} for this number.\\
Do you think your approximation is too high or too low? Why?

\vfill
\vfill

\item How does $A'(1960)$ compare to $A'(1999)$? Is that good or bad?
\vfill
\end{enumerate}


%%%%%%%%%%%%%%%%%%%%%%%%%%%%%%%%%%%%%%%%%%%%%%%%%%%%%%%%%
\pagebreak
%%%%%%%%%%%%%%%%%%%%%%%%%%%%%%%%%%%%%%%%%%%%%%%%%%%%%%%%%

\section{Learning target DFa, version 3}

Correctly interpret the meaning of the derivative in context. Assign the correct units to the value of the derivative.

\hrulefill

Let $f(t)$ be the number of centimeters of rain that have fallen by time $t$ hours after midnight.
\begin{enumerate}[(a)]

% At 3am, 1.5cm of rain had fallen
\item Suppose $f(3) = 1.5$. What are the units on the 3? What are the units on the 1.5?
\vfill
\item Write a sentence explaining what $f(3) = 1.5$ means in the story.
\vfill

% At 3am, rain fell at a rate of 0.2 cm/hr
\item Suppose $f'(3) = 0.2$. What are the units on the 3? What are the units on the 0.2?
\vfill
\item Write a sentence explaining what $f'(3) = 0.2$ means in the story.\\
\textbf{Don't say ``per'' and don't say ``rate''.}

\vfill
\item Would the statement $f'(5) = -0.7$ make sense?
Why or why not?

\end{enumerate}

%%%%%%%%%%%%%%%%%%%%%%%%%%%%%%%%%%%%%%%%%%%%%%%%%%%%%%%%%
\pagebreak
%%%%%%%%%%%%%%%%%%%%%%%%%%%%%%%%%%%%%%%%%%%%%%%%%%%%%%%%%

\section{Learning target DFa, version 4}

Correctly interpret the meaning of the derivative in context. Assign the correct units to the value of the derivative.

\hrulefill

In order to prepare for the upcoming ski season, a local resort is making snow on some of its lower trails. The snow depth, in inches, at time $t$ is $D(t)$, where $t$ is the number of days we are into December (so, for instance, $t=9$ is December 9). 

\begin{enumerate}[(a)]
\item Suppose $D(15)=28$. What are the units on the 15? What are the units on the 28?
\vfill
\item Write a sentence explaining what $D(15)=28$ means in the story.
\vfill

\item Suppose $D'(15) = 1.4$. What are the units on the 15? What are the units on the 1.4?
\vfill
\item Write a sentence explaining what $D'(15) = 1.4$ means in the story.\\
\textbf{Don't say ``per'' and don't say ``rate''.}

\vfill
\end{enumerate}

\end{document}